\section{한계}
\label{sec:limitations}

첫째, Concise, Balanced, Comprehensive는 길이만 다른 target이
아니다. information selection, expected outcome coverage,
information ordering과 총 loss-token budget이 함께 달라지는 실제
verbalization recipe다. 따라서 결과는 length-only 인과 효과가
아니라 supervision completeness와 비용이 결합된 효과로 해석한다.
둘째, 개별 sample에서 Concise evidence 집합이 Balanced의
부분집합이 되도록 강제하지 않는다. 세 조건은 평균 coverage가
증가하지만 strict dose--response 설계는 아니다. 셋째, 세 조건이
공통으로 유효한 화면 교집합을 사용하므로 전체 구축 corpus와 통제
실험의 규모가 다르다.

넷째, 모든 합성 설명은 하나의 teacher와 Stage 1 evidence에
의존한다. Stage 1에서 누락되거나 잘못 추출된 정보는 세 조건에
공통으로 전파될 수 있으며, Stage 1 coverage는 screenshot 전체의
완전한 ground truth가 아니다. 자동 coverage와 unsupported-claim
판정도 사람 검수 표본과 일치도를 보고해야 한다. 다섯째, 학습은
단일 backbone과 제한된 seed로 수행되어 run 간 분산을 충분히
추정하지 못한다.

공통 UGround 데이터에는 명목 좌표 범위를 벗어난 예제가 일부
존재하며, trajectory의 이미지 pixel 상한도 PT와 다르다. 이는
조건 간에는 공유되지만 절대 성능에는 영향을 줄 수 있다. GUI
grounding 정확도는 target 크기, 입력 해상도와 instruction 표현에
민감하므로 \citep{li2025screenspotpro,liu2025uie2i,jandial2026grounders},
공개 benchmark의 종합 점수만으로 모든 실제 UI 조건에 일반화할 수
없다. 마지막으로 약 2.6M 예제를 사용하는 영어ㆍ한국어
Comprehensive recipe는 통제 실험과 규모와 언어가 다르며, 이 결과는
실용적 효용으로만 보고한다.
